\documentclass[11pt]{article}

\usepackage[margin=1in]{geometry}
\usepackage{amsmath,amssymb,amsthm}
\usepackage{booktabs}
\usepackage{enumitem}
\usepackage{hyperref}

\newtheorem{proposition}{Proposition}
\newtheorem{definition}{Definition}
\newtheorem{remark}{Remark}

\newcommand{\R}{\mathbb R}
\newcommand{\dd}{\,\mathrm d}
\newcommand{\calZ}{\mathcal Z}
\newcommand{\calH}{\mathcal H}
\newcommand{\calB}{\mathcal B}

\title{Finite-Support-Certified Fixed Fits for the Zhao--Cui Fixed-Branch Adaptation}
\author{BayesFilter P72 planning note}
\date{2026-06-17}

\begin{document}
\maketitle

\begin{abstract}
This note explains the P70 Phase 6h failure and the proposed repair before any
new implementation work.  The failure is not that the Zhao--Cui adaptive
algorithm is wrong.  The failure is in our fixed-branch specialization: a
deterministic least-squares square-root tensor-train fit can have a very small
residual on its fitting cloud while growing by many orders of magnitude on
nearby diagnostic or replay points.  The repair proposed here is a
finite-support-certified fixed fit.  It augments the fixed least-squares
construction with predeclared guard sets, regularization and conditioning
gates, and holdout/replay admission tests.  The certificate is finite and
diagnostic: it concerns declared clouds and declared line probes, not the whole
continuum support.  This is a fixed-HMC stabilization layer unless an
individual operation is later tied to Zhao--Cui paper and author-source
anchors.
\end{abstract}

\section{Scope And Nonclaims}

The active evidence is the P70 Phase 6h result.  It reports two blockers:
\begin{enumerate}[leftmargin=2em]
    \item The degree-one, rank-two row has small fit residuals but enormous
    holdout and replay residuals.  Line probes from fitting points toward
    diagnostic points reach approximately \(2.053\times 10^{10}\) at step 1
    and \(8.987\times 10^7\) at step 2.
    \item The rank-three row fails a conditioning gate.  The failing scaled
    augmented system has smallest singular value approximately
    \(1.925\times 10^{-14}\), largest singular value approximately \(2.512\),
    and scaled condition number approximately \(1.305\times 10^{14}\).
\end{enumerate}
The raw holdout and replay residuals are already enormous, so the first
blocker is not merely a normalization artifact.  Phase 7 therefore remains
blocked.

This note does not claim that the original Zhao--Cui method fails.  Zhao and
Cui use an adaptive tensor-train route with adaptive rank, pivots, samples,
transport maps, and fitting choices.  Our branch is different: it freezes the
discrete approximation choices so that a fixed scalar can be differentiated.
Every repair in this note must therefore be classified as one of:
\[
    \text{source\_faithful},\qquad
    \text{fixed\_hmc\_adaptation},\qquad
    \text{extension\_or\_invention}.
\]
An operation is not source-faithful unless it has both paper anchors and local
author-source file/line anchors.  The phrase ``same scalar'' below means the
same scalar-valued likelihood approximation after the branch has been frozen;
it is not a claim that the adaptive Zhao--Cui computation is differentiable as
a whole.

\section{The Fixed Square-Root Fit}

Fix a time step \(t\), an external parameter value \(\beta\), and a frozen
branch
\[
    \calB_t
    =
    (\Omega_t,\Psi_t,\lambda_t,\tau_t,c_t,
      \text{basis},\text{rank},\text{sweep schedule},\text{fitting data}).
\]
The fixed branch defines transformed coordinates \(z\in\Omega_t\), a
reference density \(\lambda_t\), and a shifted square-root target
\begin{equation}
    y_t(z;\beta)
    =
    \exp\!\left\{-\frac{1}{2}
    \bigl(\Phi_t(z;\beta)-c_t\bigr)\right\}.
    \label{eq:p72-target}
\end{equation}
The fitted square-root tensor train is an element
\[
    h_\theta \in \calH_{\calB_t},
\]
where \(\theta\) denotes all TT core coefficients under the frozen basis and
rank.  The squared density used by the filter has the form
\begin{equation}
    \widehat q_t(z;\beta)
    =
    h_\theta(z;\beta)^2+\tau_t\lambda_t(z).
    \label{eq:p72-squared-density}
\end{equation}
The current fixed least-squares form fits \(h_\theta\) on a finite cloud
\(\calZ_{\rm fit}=\{z_i\}_{i=1}^{n_f}\) with weights \(w_i>0\):
\begin{equation}
    \min_\theta
    \sum_{i=1}^{n_f} w_i
    \bigl(h_\theta(z_i)-y_t(z_i)\bigr)^2
    +
    \lambda\|\Gamma\theta\|_2^2 .
    \label{eq:p72-current-objective}
\end{equation}
Here \(\Gamma\) is a coefficient penalty; \(\Gamma=I\) gives ordinary ridge.
In a single alternating least-squares core update, all other cores are fixed,
and the local problem has the linear form
\begin{equation}
    \min_{\theta_k}
    \|W_f^{1/2}(A_f\theta_k-y_f)\|_2^2
    +
    \lambda\|\Gamma_k\theta_k\|_2^2 .
    \label{eq:p72-local-ls}
\end{equation}

\section{Why Fit Residual Does Not Certify Support Stability}

\begin{proposition}[Fit residual does not control off-cloud values]
\label{prop:p72-fit-not-support}
Let \(A_f\in\R^{n_f\times m}\) be the local least-squares design matrix for a
fixed core update on the fitting cloud, and let \(a(u)^\top\in\R^{1\times m}\)
be the design row for a point \(u\notin\calZ_{\rm fit}\).  Suppose there exists
\(\delta\in\ker(A_f)\) with \(a(u)^\top\delta\ne0\).  Then, for any local
coefficient vector \(\theta\) and any scalar \(M>0\), there exists another
coefficient vector \(\theta'\) whose fitting residual is identical to that of
\(\theta\) but whose value at \(u\) differs by at least \(M\).
\end{proposition}

\begin{proof}
Set \(\theta'=\theta+s\delta\).  Since \(A_f\delta=0\),
\[
    A_f\theta'-y_f=A_f\theta-y_f.
\]
Thus the fitting residual on \(\calZ_{\rm fit}\) is unchanged.  At \(u\),
\[
    a(u)^\top\theta'-a(u)^\top\theta
    =
    s\,a(u)^\top\delta .
\]
Because \(a(u)^\top\delta\ne0\), choosing
\(|s|\ge M/|a(u)^\top\delta|\) gives the stated discrepancy.  The same
argument applies to near-null directions with small but nonzero
\(\|A_f\delta\|\): the fit residual changes little while the value at \(u\)
can change substantially.
\end{proof}

\begin{remark}
The proposition is deliberately elementary.  It is the finite-dimensional
reason the Phase 6h observation is possible.  In a TT alternating fit, each
core update is linear after the neighboring cores are frozen.  Therefore a
small empirical residual for the local design matrix cannot by itself certify
the square-root train on replay, holdout, or integration support.  The
proposition is about the empirical fit residual, not the full regularized
objective.  A sufficiently strong penalty can suppress some null directions,
but the Phase 6h failure shows that the present regularization and gates did
not suppress the observed off-cloud growth.
\end{remark}

\section{Support-Certified Fixed Fit}

The repair is to distinguish three finite sets:
\[
    \calZ_{\rm fit},\qquad
    \calZ_{\rm guard},\qquad
    \calZ_{\rm audit}.
\]
The fitting cloud \(\calZ_{\rm fit}\) defines the basic regression problem.
The guard cloud \(\calZ_{\rm guard}\) is predeclared before the fit is
observed and is meant to cover support regions the filter will later use:
proposal points, replay-like points, collocation points, and deterministic
line-probe points generated from the frozen branch.  The audit cloud
\(\calZ_{\rm audit}\) is not used to choose coefficients; it is used to decide
whether the repaired branch may advance.  Unless a phase explicitly declares
and reviews a surrogate, the guard target is the same frozen-branch target
\[
    \widetilde y_j = y_t(u_j;\beta)
\]
from \eqref{eq:p72-target}, evaluated at \(u_j\in\calZ_{\rm guard}\).

\begin{definition}[Finite-support-certified fixed objective]
\label{def:p72-support-certified-objective}
A finite-support-certified fixed fit minimizes
\begin{equation}
\begin{aligned}
    J(\theta)
    &=
    \sum_{z_i\in\calZ_{\rm fit}}
    w_i\bigl(h_\theta(z_i)-y_i\bigr)^2
    +
    \alpha
    \sum_{u_j\in\calZ_{\rm guard}}
    \widetilde w_j\bigl(h_\theta(u_j)-\widetilde y_j\bigr)^2  \\
    &\quad
    +
    \lambda\|\Gamma\theta\|_2^2
    +
    \rho\,\mathcal R_{\rm shape}(\theta),
\end{aligned}
    \label{eq:p72-support-certified-objective}
\end{equation}
where \(\alpha,\lambda,\rho\ge0\) are fixed before execution.  The optional
shape penalty \(\mathcal R_{\rm shape}\) may penalize derivative energy,
endpoint growth, or line-probe growth, but only after being explicitly
classified as a fixed-HMC adaptation or an extension.
\end{definition}

\begin{proposition}[Finite guard control]
\label{prop:p72-guard-control}
Suppose \(J(\theta)\le \varepsilon^2\) in
\eqref{eq:p72-support-certified-objective} with nonnegative terms and
\(\alpha>0\).  Then
\[
    \sum_{z_i\in\calZ_{\rm fit}}
    w_i\bigl(h_\theta(z_i)-y_i\bigr)^2
    \le \varepsilon^2,
\]
and
\[
    \sum_{u_j\in\calZ_{\rm guard}}
    \widetilde w_j\bigl(h_\theta(u_j)-\widetilde y_j\bigr)^2
    \le \varepsilon^2/\alpha .
\]
If \(\lambda>0\), then
\[
    \|\Gamma\theta\|_2^2\le \varepsilon^2/\lambda .
\]
\end{proposition}

\begin{proof}
Each displayed quantity is a nonnegative summand of \(J(\theta)\), up to its
positive multiplier.  Dropping all other nonnegative terms gives the three
inequalities.
\end{proof}

This proposition is not a theorem about the continuum target.  It says
something narrower but operationally important: if a finite guard set is part
of the objective, then small objective value controls the weighted aggregate
fit on that guard set.  It does not by itself rule out a large pointwise spike
at a lightly weighted point or away from the guard set.  For that reason the
admission gate below includes explicit finite maximum checks.  Any continuum
support claim still requires design-measure coverage, additional analysis, or
empirical validation.

\section{Conditioning Gate For Local ALS Solves}

The row-B failure is a conditioning failure.  For each local core update, let
\[
    \bar A_k
    =
    \begin{bmatrix}
      W^{1/2} A_k D_k^{-1}\\
      \sqrt{\lambda}\,\Gamma_k D_k^{-1}
    \end{bmatrix},
    \label{eq:p72-scaled-augmented-design}
\]
where \(D_k\) is a deterministic invertible diagonal column scaling matrix,
fixed before the solve is inspected.  The local update is admissible only if
the singular values of \(\bar A_k\) satisfy predeclared thresholds:
\begin{equation}
    \sigma_{\min}(\bar A_k)\ge \sigma_{\min}^{\rm gate},
    \qquad
    \kappa(\bar A_k)
    =
    \frac{\sigma_{\max}(\bar A_k)}{\sigma_{\min}(\bar A_k)}
    \le \kappa_{\max},
    \label{eq:p72-conditioning-gate}
\end{equation}
and the effective rank is not below a predeclared floor.  The effective-rank
convention must also be fixed before execution, for example by declaring a
relative singular-value cutoff and whether the ratio is taken against the
number of columns or the declared local degrees of freedom.  A rank increase is
therefore not automatically an improvement.  If the added rank directions
create near-null directions, the fit is rejected or repaired before any
validation claim.

\section{Admission Gate}

For a line probe, choose a fitting point \(z\), a guard or audit point \(u\),
and a predeclared set of fractions \(S\subset[0,1]\).  The finite probe path is
\[
    \gamma_{z,u}(s)=(1-s)z+su,\qquad s\in S.
\]
The admission record must predeclare whether it uses an absolute value ceiling,
a residual ceiling, a growth ratio relative to fit-scale, or an overshoot ratio
relative to endpoints.  A typical absolute finite check has the form
\[
    \max_{(z,u),\,s\in S}|h_\theta(\gamma_{z,u}(s))|\le G_{\max},
\]
while a residual check has the form
\[
    \max_{v\in\calZ_{\rm guard}\cup\calZ_{\rm audit}}
    |h_\theta(v)-y_t(v;\beta)|\le R_{\max}.
\]

After fitting, a branch may be admitted only if all of the following
predeclared checks pass:
\begin{enumerate}[leftmargin=2em]
    \item fit residual on \(\calZ_{\rm fit}\);
    \item guard residual on \(\calZ_{\rm guard}\);
    \item aggregate and maximum guard residuals on \(\calZ_{\rm guard}\);
    \item aggregate and maximum holdout/replay residuals on \(\calZ_{\rm audit}\);
    \item declared line-probe value, residual, growth, or overshoot checks from
    fitting points toward guard or audit points;
    \item squared-density normalizer and defensive-mass checks;
    \item local design conditioning checks from \eqref{eq:p72-conditioning-gate};
    \item checks that declared rank directions are actually used, so aggregate
    residuals cannot hide inactive declared directions.
\end{enumerate}
The in-sample residual alone is never a promotion criterion.

\begin{proposition}[Admission excludes the observed Phase 6h failure pattern]
\label{prop:p72-admission-excludes-phase6h}
Assume the admission gate includes a declared finite line-probe ceiling
\(G_{\max}\), a maximum raw holdout/replay residual ceiling \(R_{\max}\), and
the conditioning gate \eqref{eq:p72-conditioning-gate}.  Then any admitted
branch has no measured finite line-probe statistic above \(G_{\max}\), no
measured raw holdout/replay residual above \(R_{\max}\), and no admitted local
ALS system with condition number above \(\kappa_{\max}\).
\end{proposition}

\begin{proof}
This is immediate from the definition of admission.  The point is not depth
but discipline: the quantities that failed in Phase 6h become veto diagnostics
instead of after-the-fact explanations.
\end{proof}

\section{Source-Governance Classification}

\begin{center}
\begin{tabular}{p{0.25\linewidth}p{0.29\linewidth}p{0.32\linewidth}}
\toprule
Operation & Classification & Reason \\
\midrule
Squared TT/SIRT density, defensive density, sequential normalizer role &
source\_faithful when cited with Zhao--Cui paper and author-source anchors &
These are part of the inspected Zhao--Cui route and existing local ledgers. \\
Freezing ranks, bases, domains, samples, sweep schedule, and branch identity &
fixed\_hmc\_adaptation &
This preserves the broad route while making a differentiable same-scalar
branch.  It is not the adaptive author algorithm. \\
UKF-guided center, scale, or covariance orientation &
fixed\_hmc\_adaptation unless separately source-anchored &
UKF is a scout for branch design, not a correctness oracle. \\
Guard/collocation objective terms, line-growth penalties, and explicit
off-cloud admission gates &
extension\_or\_invention unless tied to a source route in a later audit &
These directly address the fixed-variant stability defect.  They cannot close
a Zhao--Cui source-faithfulness gap by themselves. \\
Christoffel, leverage-score, or stable-sampling rules for least squares &
extension\_or\_invention until the relevant literature is audited &
They are plausible numerical-stability tools, but must be separately sourced
before being presented as mathematical support. \\
\bottomrule
\end{tabular}
\end{center}

\section{Literature Directions}

The local bibliography already contains the direct TT and filtering anchors:
Zhao and Cui's sequential TT state-space method, Oseledets' tensor-train
decomposition, Oseledets--Tyrtyshnikov TT-cross approximation, and
Cui--Dolgov squared inverse Rosenblatt transports.  These support the broad
TT/SIRT context when the precise inspected sections are cited.

The proposed stabilization also points toward the literature on stable
weighted least squares, oversampling, Christoffel or leverage-score sampling,
and polynomial approximation.  Those papers should be treated as candidate
support until their technical sections are audited.  They may motivate a
support-certified design, but they do not yet justify a theorem-level claim in
the BayesFilter monograph.

\section{Implementation Consequences}

The next implementation phase should not go to Phase 7.  It should design and
implement a Phase 6i repair with the following minimal contract:
\begin{enumerate}[leftmargin=2em]
    \item define predeclared guard and audit clouds from the frozen branch;
    \item implement a support-certified fixed objective or a guard-gated
    objective-preserving variant;
    \item preserve TensorFlow/TensorFlow Probability as the BayesFilter
    algorithmic backend;
    \item record the column scaling convention, singular spectra, effective
    rank convention, and conditioning for every admitted local solve;
    \item reject a branch if off-cloud growth, holdout/replay residuals,
    normalizer checks, or conditioning gates fail;
    \item label every new operation as fixed-HMC adaptation or extension unless
    source anchors prove source faithfulness.
\end{enumerate}

\section{What Would Count As Success}

A repaired branch has not succeeded because it fits \(\calZ_{\rm fit}\).  It
has only passed the lower repair gate if, on the same bounded diagnostic scale
as Phase 6h, it shows all of the following:
\[
\begin{gathered}
    \text{finite normalizer},\qquad
    \text{bounded guard residual},\qquad
    \text{bounded holdout/replay residual},\\
    \text{bounded line growth},\qquad
    \text{admissible conditioning}.
\end{gathered}
\]
Even then, the conclusion is only that the fixed-variant lower gate has been
repaired enough to consider the next reviewed diagnostic.  It is not a d18
accuracy claim, not a scaling claim, not HMC readiness, and not adaptive
Zhao--Cui parity.

\begin{thebibliography}{9}

\bibitem{zhao-cui-2024}
Yiran Zhao and Tiangang Cui.
\newblock Tensor-Train Methods for Sequential State and Parameter Learning in
State-Space Models.
\newblock \emph{Journal of Machine Learning Research}, 25:1--51, 2024.

\bibitem{oseledets-2011}
Ivan V. Oseledets.
\newblock Tensor-Train Decomposition.
\newblock \emph{SIAM Journal on Scientific Computing}, 33(5):2295--2317, 2011.

\bibitem{oseledets-tyrtyshnikov-2010}
Ivan Oseledets and Eugene Tyrtyshnikov.
\newblock TT-Cross Approximation for Multidimensional Arrays.
\newblock \emph{Linear Algebra and its Applications}, 432:70--88, 2010.

\bibitem{cui-dolgov-2021}
Tiangang Cui and Sergey Dolgov.
\newblock Deep Composition of Tensor-Trains Using Squared Inverse Rosenblatt
Transports.
\newblock arXiv:2007.06968, 2021.

\end{thebibliography}

\end{document}
